\documentclass[11pt]{article}
\usepackage[utf8]{inputenc}\usepackage[T1]{fontenc}\usepackage{lmodern}
\usepackage[margin=1in]{geometry}
\usepackage{microtype}\setlength{\emergencystretch}{3em}
\usepackage{amsmath,amssymb,booktabs,graphicx,array,xcolor,enumitem,caption}
\graphicspath{{figures/}}
\usepackage[round]{natbib}
\usepackage[hidelinks]{hyperref}



\title{\bf Decomposed Scaling Forecasts Need a Derived Total, Not a Fitted One}


  \author{\normalsize\itshape Author names and affiliations removed for
  double-anonymous peer review}

\date{}

\begin{document}
\maketitle

\begin{abstract}
Organisations increasingly decide where to spend an artificial intelligence budget
by asking which of a system's failures the next model generation will remove on its
own. The usual evidence is a decomposition: sort the errors into categories, fit a
scaling curve to each, and read off which shrink fastest. Studies that do this
commonly also fit the aggregate error rate as a curve of its own. We show the two
cannot both be right. If the categories follow power laws with different exponents,
their sum is provably not a power law, so an independently fitted aggregate
contradicts the categories it summarises, and the contradiction grows without
bound. On our measurements the parts reach 1.27 times the fitted
whole at one billion parameters and 2.82 times at one trillion,
exceeding ten percent beyond 174M parameters
[57M, 635M]. The repair is to derive the
total from the parts rather than fit it separately. What that costs is the
aggregate exponent itself, which does not exist: its effective slope drifts from
$-$0.299 to $-$0.174, covering 27\% of the
span of the category exponents it is meant to summarise. We contribute a diagnostic
needing only the coefficients a study already prints, an audit finding that speech
recognition already follows the derived-total convention while scaling analyses do
not, and the observation that conclusions resting on such a pair's out-of-range
behaviour are void, which we show by withdrawing one of our own.
\end{abstract}

\vspace{0.5em}
\noindent\textbf{Keywords:} scaling laws; error decomposition; design science;
model evaluation; artificial intelligence governance; forecast admissibility

% ====================================================================
\section{Introduction}
\label{sec:intro}

A firm running a document-processing system, a coding assistant, or a clinical
triage aid faces a recurring question with a large budget attached. The system
fails in several distinguishable ways. Some of those failures the next model
generation will probably remove without anyone doing anything. Others it will not.
Money spent mitigating the first kind is wasted, and money withheld from the second
kind buys a defect that persists through every upgrade. Deciding which is which is
a forecasting problem, and the field has a standard way of answering it: partition
the observed errors into categories, fit a scaling curve to each category, and
compare the exponents.

Studies that do this almost always report one more thing: a scaling curve for the
aggregate error rate, fitted separately. That aggregate is usually the headline. We
show the two reports are mutually inconsistent, and that the inconsistency is forced
rather than incidental. If the categories are power laws with different exponents,
their sum is provably not a power law, because a sum of power laws has a slope that
moves. An aggregate fitted as a power law therefore contradicts the categories it is
meant to summarise, and the contradiction grows without bound with scale. Past a
computable point the fitted parts exceed the fitted whole, which is not a wide
confidence interval but an arithmetic impossibility.

The repair is simpler than the diagnosis, and it is not ours. Stop fitting the
aggregate and define it as the sum of the categories. That is consistent at every
scale, preserves every category exponent, and needs no new machinery. What it costs
is the aggregate exponent, which under the repair does not exist as a single number:
on our data its effective slope drifts from $-$0.299 at the bottom of
our range to $-$0.174 at one billion parameters, covering
27\% of the span of the category exponents it was meant to
summarise. A study may report per-category exponents, or an aggregate exponent, but
reporting both is reporting two things that cannot both hold.

A companion study of ours
 (Anonymous, 2026){} decomposed the illegal
moves made by chess transformers, fitted a curve per category and a separate curve
for the aggregate, and projected to one billion parameters. The inconsistency
exceeded ten percent beyond 174M parameters
[57M, 635M], well below where we read it. The
fits had passed every quality gate we had and agreed closely with the data in range.
We report the diagnosis, the repair, and our own correction together, because the
sequence is the argument: this failure is invisible to the checks practitioners
actually run.

\paragraph{Why this matters beyond the laboratory.}
The decomposition-and-extrapolate pattern has escaped the research literature. It
appears in vendor roadmaps (``this class of error is on a steep curve''), in
internal capability planning, in procurement comparisons, and in assurance
arguments submitted to auditors. A manager reading such a claim currently has no
way to check it other than trusting the analyst, because the defect is not visible
in any of the artefacts normally shown: the fit looks excellent, the residuals are
small, the intervals are narrow. The diagnostic we provide changes that. It needs
only the exponents a study already prints, it runs in milliseconds, and it returns
a scale past which the forecast is not entitled to its own numbers. That converts a
question of analytic authority into a question anyone can audit.

\paragraph{Research questions and objectives.}
\begin{description}[nosep,leftmargin=2.2em]
\item[RQ1] Is the inconsistency between per-category curves and a separately fitted
aggregate incidental to particular data, or forced?
\item[RO1] Establish the condition under which the fitted parts must outgrow the
fitted whole, and derive the scale at which a given pair becomes inconsistent
(Section~\ref{sec:design}).
\item[RQ2] What does the derived-total repair cost, and is further machinery needed?
\item[RO2] Compare the derived total against the fitted one, and compare the two
admissible ways of writing the composition under a common loss
(Sections~\ref{sec:design} and~\ref{sec:eval}).
\item[RQ3] Can the defect be detected from what published studies already print,
without access to their data?
\item[RO3] Implement a diagnostic consuming only reported coefficients and apply it
to published decompositions (Sections~\ref{sec:arch} and~\ref{sec:eval}).
\item[RQ4] What does an inconsistent pair invalidate, beyond the forecast itself?
\item[RO4] Test whether a conclusion we had drawn against such a pair survives being
redrawn admissibly (Section~\ref{sec:eval}).
\end{description}

\paragraph{Contributions.}
The contribution is a diagnostic artifact together with the design knowledge it
carries, in three pillars. First, a consistency condition for the pair of a
per-category decomposition and a fitted aggregate, with a computable breakdown
scale, showing the conflict is forced and identifying the aggregate as the term that
must give way. Second, a characterisation of what the repair costs, namely the
aggregate exponent, quantified as the drift of the derived total's slope, together
with the finding that the two admissible ways of writing the composition agree to
within 
esult{admMaxShareGap} in share and so present no further choice worth
arguing about. Third, evidence that a conclusion resting on such a pair's
out-of-range behaviour is void, demonstrated by withdrawing a rejection of our own
that turned out to rest on nothing.

% ====================================================================
\section{Related Work}
\label{sec:related}

\paragraph{Scaling laws and what is fitted to them.}
Power-law relationships between model scale and loss \citep{kaplan2020scaling} and
their compute-optimal refinement \citep{hoffmann2022training} established the
practice of extrapolating fitted curves to scales not yet trained. Those studies fit
a single aggregate quantity, for which no partition identity exists and the
criticism here does not apply. The difficulty arrives when the same machinery is
applied category by category to a decomposed error, which is what practitioners
want because it is category-level behaviour that determines where to spend. We
differ from this literature by asking not how accurate such fits are but whether
they are admissible at all.

\paragraph{Metric artefacts in scale measurement.}
\citet{schaeffer2023emergent} showed that apparently discontinuous capability gains
can be manufactured by the choice of metric, and that linear or continuous metrics
produce smooth curves from the same model outputs. That result and ours are the
same genre of criticism, that a conclusion about models turned out to be a
conclusion about the measurement apparatus, and they are complementary in
mechanism: theirs concerns the shape imposed by a discontinuous scoring function,
ours the constraint violated by fitting the parts of a whole independently. A study
can satisfy their recommendation completely, using a perfectly continuous metric,
and still produce the inconsistent pair we describe.

\paragraph{Compositional data and non-power-law scaling.}
Modelling a vector of shares with an additive-logratio link is the standard
apparatus of compositional data analysis \citep{aitchison1986statistical}, and
Equation~\ref{eq:softmax} is that apparatus with $\log N$ as the covariate. We claim
no novelty for the model; the contribution is the observation that a decomposition
plus a fitted aggregate is already an implicit compositional model, and an
inconsistent one, so the field has been doing compositional analysis without the
constraint that defines it. Separately, work on scaling laws with more than one
regime \citep{caballero2023broken} arrives at the same conclusion that a single
power law is often the wrong form, from the shape of observed curves rather than
from a constraint the curves must jointly satisfy. Our argument is complementary and
stronger in one narrow respect: it says a single power law for the aggregate cannot
be right, without needing to observe the aggregate at all.

\paragraph{Design science and the evaluated artifact.}
We frame the contribution as design-science research in the sense of
\citet{hevner2004design}, following the process of \citet{peffers2007design} and
positioning the output as an improvement-type contribution in the terms of
\citet{gregor2013positioning}: a known problem, forecasting where scale will help,
addressed by a new solution whose transferable content is the admissibility
condition rather than the code that computes it. The artifact is evaluated against
measurements from two independent model families, one of which we did not produce.

\paragraph{The instrument domain.}
Our measurements come from chess transformers, where a rules engine supplies an
exact and cheap oracle for whether an output is legal and for which constraint it
violated \citep{karvonen2024emergent}. That is why the domain is used here: it
yields an error decomposition that is unambiguous and not itself a modelling
choice. Nothing in the diagnostic is specific to it, and
Section~\ref{sec:implications} sets out where else the pattern applies.

% ====================================================================
\section{Design Problem and Requirements}
\label{sec:design}

\paragraph{The problem, stated plainly.}
A study reports, for each of $k$ error categories, a fitted relationship
$r_i(N) = A_i N^{a_i}$ between model scale $N$ and the category's rate, and a
separate fit $T(N) = B N^{b}$ for the total error rate. A reader wants to know what
the composition looks like at some $N$ beyond the fitted range. The categories
partition the errors, so the identity
\begin{equation}
\sum_{i=1}^{k} r_i(N) \;=\; T(N) \qquad \text{for all } N
\label{eq:closure}
\end{equation}
holds of the quantities being estimated. It does not hold of the estimates.

\paragraph{The consistency condition.}
Let $a^\star = \max_i a_i$. The sum of the parts is asymptotically proportional to
$N^{a^\star}$ and the fitted whole to $N^{b}$, so the ratio of parts to whole grows
as $N^{a^\star - b}$ and diverges precisely when $a^\star > b$. That inequality is
not a coincidence to be checked case by case. Fitted within the observed range, $b$
tracks the rate-weighted mean of the $a_i$, because the total is the rate-weighted
sum of the parts there. A maximum is at least a weighted mean, with equality only
if every $a_i$ is identical. Hence:

\begin{quote}
\emph{If the categories are fitted as power laws with exponents that are not all
equal, no power law can represent their total, and any separately fitted aggregate
must contradict them beyond a finite scale computable from the fitted coefficients.}
\end{quote}

Since exponents that are all equal would make the decomposition pointless, every
decomposition worth reporting conflicts with its own fitted aggregate somewhere. Our
values illustrate it: $a^\star$ is $-$0.112, the rate-weighted mean of the
part exponents is $-$0.266, and the separately fitted $b$ is
$-$0.267.

\paragraph{Which term is wrong, and the obvious repair.}
It is worth being precise about where the error lies, because the natural first
reading is that the category fits are at fault and they are not. A sum of power laws
with distinct exponents is not a power law. The categories may each be perfectly
well specified; what cannot be sustained is the additional claim that their total is
also a power law. The misspecified object is the aggregate.

The repair follows immediately and requires nothing new: do not fit the aggregate,
and define
\begin{equation}
T(N) \;:=\; \sum_{i=1}^{k} r_i(N),
\label{eq:derived}
\end{equation}
which satisfies Equation~\ref{eq:closure} identically at every scale, for every
setting of the coefficients, and preserves each category exponent exactly. No link
function, no constrained estimation, and no reparameterisation is needed. We state
this plainly because it is the correct remedy and because an earlier version of our
own analysis reached for a more elaborate one before noticing it.

\paragraph{What the repair costs.}
It costs the aggregate exponent, and that is not a small thing, because the
aggregate exponent is usually the number a study is read for. Under
Equation~\ref{eq:derived} the total has a slope that moves with scale. On our
measurements the effective slope $\mathrm{d}\log T/\mathrm{d}\log N$ is
$-$0.299 at the bottom of the observed range, $-$0.228 at
the top, $-$0.174 at one billion parameters and $-$0.137 at one
trillion. The drift within the observed range alone is 0.071, and
by one billion it is 0.125, which is 27\% of the
0.461 span between our steepest and shallowest category exponents. An
aggregate exponent quoted from such a fit is a local slope valid over the fitted
range, and reporting it as a property of the system is the substantive error behind
the arithmetic one. This is the same observation that motivates work on scaling laws
with more than one regime \citep{caballero2023broken}, arrived at from the
constraint side rather than the curve-shape side.

\paragraph{Writing the composition.}
Dividing Equation~\ref{eq:derived} through gives the category shares
\begin{equation}
s_i(N) \;=\; \frac{A_i N^{a_i}}{\sum_j A_j N^{a_j}}
      \;=\; \frac{\exp(\log A_i + a_i \log N)}{\sum_j \exp(\log A_j + a_j \log N)},
\label{eq:softmax}
\end{equation}
which is a softmax linear in $\log N$, that is, a compositional model with an
additive-logratio link and $\log N$ as the covariate, in the sense of
\citet{aitchison1986statistical}. So the derived-total repair and a compositional
fit on the simplex are the same model of the composition; they differ only in the
loss used to estimate it, log-rate space in the first case and share space in the
second. Section~\ref{sec:eval} reports that on our data they agree to within
0.0060 in share across six orders of magnitude and agree on the
asymptotics, so an analyst need not choose carefully between them. What matters is
that either is used instead of the fitted aggregate.

\paragraph{Design principles.}
Table~\ref{tab:dp} states the principles the artifact realises. They are stated so
as to transfer: none mentions chess, error legality, or any property of our
instrument.

\begin{table}[t]
\centering
\caption{Design principles and their realisation in the diagnostic.}
\label{tab:dp}
\small
\begin{tabular}{@{}p{0.30\linewidth}p{0.62\linewidth}@{}}
\toprule
\textbf{Design principle} & \textbf{Realisation} \\
\midrule
DP1. Consume only what is published & The check takes exponents and intercepts, the
figures a study already prints, so it can be applied to work whose data are
unavailable. \\
DP2. Separate the verdict from the remedy & The diagnostic reports admissibility
and a breakdown scale without requiring the analyst to adopt any particular
replacement model. \\
DP3. Report a scale, not a verdict alone & A decomposition is not simply good or
bad; it is usable up to a threshold, so the output is the threshold. \\
DP4. Make the admissible alternative cost nothing & The closed reformulation is
chosen so that in-range accuracy does not decline, removing the incentive to skip
it. \\
DP5. Fail loudly on inputs that cannot bear a fit & A category with a zero rate at
any scale has no exponent; the implementation refuses rather than flooring the
value and inventing a slope. \\
DP6. State what is invalidated & Output distinguishes what an inconsistent pair
costs, which includes conclusions drawn against it, not only its forecasts. \\
\bottomrule
\end{tabular}
\end{table}

\paragraph{In scope and out of scope.}
In scope: decompositions in which categories partition a measured total, each fitted
as a power law in a common scale variable. Out of scope: single aggregate scaling
laws, for which no partition identity exists and the criticism does not apply;
categories that overlap rather than partition, for which a different bound is
needed; and the question of whether any particular fitted exponent is correct, which
this work does not address and does not need to. The conflict arises between fits
that may each be individually excellent, which is precisely what makes it hard to
see.

% ====================================================================
\section{Artifact Architecture}
\label{sec:arch}

The artifact has three separable parts, deliberately, so that an analyst may adopt
one without the others.

\paragraph{The checker.}
A module that takes a dictionary of category coefficients and a total, and returns
the divergence rate $a^\star - b$, an admissibility verdict, the breakdown scale at
a chosen tolerance, and the overspend ratio at any scales of interest. It depends on
one numerical library and nothing else, including nothing from the study that
produced our measurements, so that it can be pointed at published work in unrelated
fields.

\paragraph{The estimator.}
A convenience path that fits the coefficients from raw measurements rather than
requiring them to be transcribed, applying DP5: a category whose rate is zero at
some scale is refused, with an explanation, rather than silently floored.

\paragraph{The derived total.}
The remedy of Equation~\ref{eq:derived}: the total is computed from the category
curves rather than fitted, so Equation~\ref{eq:closure} holds identically and no
setting of the coefficients can produce an impossible forecast. The module also
returns the effective slope at any scale, because an analyst who previously quoted
an aggregate exponent needs to see what has replaced it, and returns the drift
across a stated range so that a single number can still be quoted where the drift
is small enough to tolerate. The equivalent compositional fit of
Equation~\ref{eq:softmax} is provided alongside it for analysts who prefer to
estimate in share space; the two are the same model and the implementation says so
rather than presenting them as alternatives to be chosen between.

\paragraph{Tolerance.}
Admissibility is asymptotic, but in-range fit noise makes a strict threshold
useless, so the breakdown scale is defined at a tolerance: the scale at which the
parts exceed the whole by more than a stated fraction, ten percent by default. The
search starts at the top of the observed range, because the ratio of parts to whole
is U-shaped, large at small scales where the steep categories dominate and large at
large scales where the shallow one does. Searching upward from zero finds the wrong
arm and reports a threshold beneath the data, which is a mistake we made and fixed.

% ====================================================================
\section{Method and Measurements}
\label{sec:method}

\paragraph{Instrument.}
Measurements come from 18 decoder-only transformers spanning
3.4M to 39.9M parameters, 6 sizes at
3 seeds each, trained on human chess games and evaluated without
search. Every model output from ply twenty onward is checked by a rules engine and,
when illegal, assigned to exactly one category by what made it illegal: no piece on
the source square, an opponent's piece there, the mover's own piece on the
destination, a geometric violation, or a move that is legal in every respect except
that it leaves the king in check. The classes are ordered by how much of the board
must be consulted to detect them, which is why their exponents differ and why the
decomposition is informative in the first place.

\paragraph{Independent replication.}
The second family is a publicly released ladder of character-level chess models
trained by another laboratory on a different corpus with different code
\citep{karvonen2024emergent}, spanning a 38-fold parameter range. Its
notation does not name a source square, so its taxonomy is coarser than ours; we do
not attempt to match the taxonomies when testing admissibility, because the
diagnostic does not require it. Each family is checked against its own partition.

\paragraph{Provenance.}
Every adapter between an external model and our measurements fails silently rather
than with an error, so each was verified and the checks recorded: a reconstructed
32-character vocabulary, move-boundary alignment verified three ways
at 4,081 sites with 0 disagreements, and a
forward pass choosing a legal move at 672 of 672
positions in a third model. No number in this manuscript is typed by hand; each
enters through a generated macro tagged with the results file it came from.

\paragraph{Computing environment.}
All runs fit on a single NVIDIA GeForce RTX 3050 Laptop GPU with 4\,GiB of device
memory, in a host with 15\,GB of physical memory visible to the
operating system, under Windows 11, using Python~3.12.10 with
PyTorch~2.5.1+cu121, NumPy~2.4.6 and SciPy~1.18.1. The
memory ceiling, rather than a judgement about what was worth measuring, is what sets
the top of our ladder. It is worth stating the consequence bluntly, because it is
the underlying reason the companion forecast went wrong: the ladder spans
12-fold in parameters, and the forecast was read
25-fold beyond its top. No fitted form is owed that much
credit, and the arithmetic failure documented here is what happens when one is
extended that far without a constraint to stop it.

% ====================================================================
\section{Illustrative Walkthrough}
\label{sec:walk}

An analyst has fitted five category curves and a total to the ladder above, and
wishes to state what the composition will look like at one billion parameters. The
fits are good: each category declines monotonically across the ladder, and the
category exponents are separated well beyond their intervals, from
$-$0.574 for the most local category to $-$0.112 for the
least, a paired differential of +0.461
[+0.259, +0.599]. Nothing in the residuals
suggests trouble.

Running the check takes the eleven fitted coefficients, five categories and the
aggregate, and returns: the largest category exponent $a^\star$ is
$-$0.112 against a fitted aggregate exponent of $-$0.267, so the
parts outgrow the whole, and the inconsistency passes ten percent at
174M parameters [57M, 635M].
The intended forecast scale is one billion, several times past that threshold, and
the parts there stand at 1.27 times the whole
[1.13, 1.57], above one in
100\% of resamples of the ladder.

The analyst has one unusable option and two usable ones. The unusable option is to
widen the intervals, which does not help, because the fitted aggregate is not
imprecise but contradicted. The first usable option is to keep both fits and report
nothing above 174M parameters. The second, and the one we
recommend, is to drop the fitted aggregate and derive the total from the categories,
which is consistent at every scale. The cost of the second is that there is no
longer an aggregate exponent to quote, only a slope that moves;
Section~\ref{sec:eval} reports how far it moves here.

% ====================================================================
\section{Evaluation}
\label{sec:eval}

\paragraph{Does the defect occur in practice, and how badly?}
Table~\ref{tab:overspend} gives the ratio of independently extrapolated parts to the
independently extrapolated whole. At the top of our observed range the ratio is
1.02, which is fit noise and unobjectionable. One order of
magnitude beyond it the parts have overtaken the whole, and three orders beyond it
they are nearly three times the total they are part of.

\begin{table}[t]
\centering
\caption{Parts-to-whole ratio under independently fitted category curves. A value
above one is not a wide estimate but an impossibility. Interval from resampling the
ladder's size rungs.}
\label{tab:overspend}
\small
\begin{tabular}{@{}lrr@{}}
\toprule
\textbf{Scale} & \textbf{Parts / whole} & \textbf{Interval} \\
\midrule
Top of observed range (39.9M) & 1.02 & n/a \\
One billion parameters & 1.27 &
[1.13, 1.57] \\
One trillion parameters & 2.82 & n/a \\
\bottomrule
\end{tabular}
\end{table}

\paragraph{Does it replicate on a family we did not train?}
Directionally yes, and that is as much as three checkpoints can support. Refitting
the independent ladder over its own 4 usable categories gives a
positive divergence rate, so the same conflict exists, and a breakdown scale of the
same order of magnitude as ours. We deliberately do not quote those figures to three
significant digits. Each is a two-parameter fit to three points, one of the five
categories had to be dropped because its rate is zero at two of the three, and the
resulting coefficients carry no usable precision. The claim the replication supports
is that the conflict is not a peculiarity of our ladder, not that its magnitude has
been estimated twice. A different tokenisation, a different taxonomy, a
different corpus and a different laboratory reproduce the defect within the same
order of magnitude, which is what Section~\ref{sec:design} predicts and what makes
this a property of the method rather than of our data. One of that family's
categories is excluded because its rate is zero at two of three checkpoints; an
exponent fitted there would be an artefact of the floor, and DP5 makes the
implementation refuse it rather than produce one.

\paragraph{Do the two admissible forms differ?}
Barely, which is the useful answer. The derived-total model and the compositional
fit are the same model estimated under different losses, so the only question is
whether the loss matters. Scored in share space the compositional fit attains
0.01251 against 0.01274 for the derived total; scored
in rate space, 0.00234 against 0.00244. The compositional
fit is marginally ahead in both, which is expected since neither is being scored in
a space that favours it over the other by construction, and the two disagree by at
most 0.0060 in share across six orders of magnitude. We draw no
conclusion from the ordering. An earlier version of this analysis reported the
compositional fit as several percent more accurate than the uncoupled baseline; that
comparison was not like for like, because the baseline had been estimated in
log-rate space and scored in share space, and we withdraw it.

The honest summary is that the repair is free in fit quality rather than beneficial.
The case for it is not accuracy but consistency, and the case against the fitted
aggregate is not that it fits badly, which it does not, but that it asserts
something the categories contradict.

\paragraph{What the repair costs, measured.}
Table~\ref{tab:drift} gives the effective slope of the derived total. Within the
observed range it moves by 0.071, and by one billion parameters it
has moved 0.125 from where it started, which is 27\%
of the 0.461 separating our steepest and shallowest category exponents.
An analyst who quoted $-$0.267 as the aggregate exponent was quoting a
local slope, and the forecast error from treating it as global is the quantity in
Table~\ref{tab:overspend}.

\begin{table}[t]
\centering
\caption{Effective slope of the derived total, $\mathrm{d}\log T/\mathrm{d}\log N$.
Constant only if all category exponents are equal; here it drifts across a
substantial fraction of the span it is meant to summarise.}
\label{tab:drift}
\small
\begin{tabular}{@{}lr@{}}
\toprule
\textbf{Scale} & \textbf{Effective slope} \\
\midrule
Bottom of observed range (3.4M) & $-$0.299 \\
Top of observed range (39.9M) & $-$0.228 \\
One billion parameters & $-$0.174 \\
One trillion parameters & $-$0.137 \\
\midrule
Separately fitted aggregate exponent & $-$0.267 \\
\bottomrule
\end{tabular}
\end{table}

\paragraph{Published decompositions: who already does this right.}
We applied the check to published error decompositions, which
Section~\ref{sec:limits} of an earlier draft had listed as undone. The binding
constraint turned out to be finding a published table that has the failing
configuration at all, and that is itself the finding.

Speech recognition is the mature case. Word error rate is defined as the sum of
substitutions, deletions and insertions, so the aggregate is derived rather than
fitted, exactly as Equation~\ref{eq:derived} recommends, and no inconsistency is
possible. Checking the identity row by row against a recent published table
\citep{loquacious2025supplementary} confirms it holds in 11 of
11 rows, with a worst discrepancy of 0.10 points, which is
what independent rounding of three components and a total to one decimal produces.
That field settled on the derived total decades ago and does not discuss it, because
within it the alternative never seemed available.

The configuration that fails is the one where the aggregate is the headline and the
categories are secondary, which is the habit of scaling analysis rather than of
error analysis. Of the 2 sources we examined in detail,
1 fits the aggregate independently, and it is ours. We report this
as a limitation converted into a finding rather than as a survey: it establishes
that the repair is already standard practice in a field that decomposes errors
carefully, and it does not establish how common the failing configuration is in the
scaling literature, which would need the survey described in FR6.

\paragraph{The same discipline outside the partition.}
A quantity confined to the unit interval needs a link that respects the confinement
even when it is not a member of a partition. The share of errors occurring despite
a correctly decoded local state, which cuts across the categories rather than
belonging to them, was observed between 0.148 and
0.321. Fitted on the logit it reaches 0.596 at one
billion and 0.951 at one trillion, remaining a proportion throughout.
Fitted linearly in the log, as the categories originally were, it would not.

\paragraph{What an inconsistent pair invalidates.}
This is the evaluation we did not expect to run. While preparing this manuscript we
tested a hypothesis that failure composition is governed by a model's achieved
competence rather than by its size, rejected it on the grounds that the fitted shares
exceeded one, and recorded the rejection. That rejection was worthless. Shares
exceeding one is what the inconsistent form produces everywhere, at every scale, for
any data whatever, so the observation carried no information about the hypothesis
under test.

We repeated the comparison admissibly, fitting the composition on the simplex in
competence coordinates, matching the two families down to the three categories both
taxonomies share, and predicting the independent family's checkpoints entirely out of
sample. The hypothesis is genuinely refuted: mean total variation distance
0.490, against 0.100 for size coordinates and
0.115 for the baseline of predicting our own mean composition
regardless of input. The conclusion survived; the original argument for it did not.

Two cautions on this result, both of which cut against us. The claim to take away is
narrow: a conclusion that depends on an inconsistent pair's out-of-range or
share-space behaviour is void, because that behaviour is determined by the
misspecification rather than by the data. A conclusion resting on in-range residuals
would have been unaffected, and we do not claim otherwise. Second, size coordinates
beat the constant baseline only by 0.015 in total variation, which
is close enough to nothing that the right reading is that neither covariate transfers
across model families, rather than that size transfers and competence does not.

\paragraph{Honest negative results.}
Two further ideas were tested and are reported because they failed. We predicted
from the closed model that error diversity should collapse with scale, since a
softmax converges to its steepest component, and measured the normalised entropy of
the composition. It declines on our ladder, with a slope whose bootstrap interval
excludes zero, and it rises on the independent ladder. One ladder is not a result,
and we do not report entropy as a scaling signature. We also computed the scale at
which the least local category overtakes the most local one, intending it as a
headline, and found it lies below our smallest model, so it describes the past
rather than the future. Neither appears in our conclusions.

% ====================================================================
\section{Discussion}
\label{sec:discussion}

\paragraph{What the integration buys.}
Each part of the artifact is individually unremarkable. A ratio of two fitted curves
is arithmetic; a softmax is standard; refusing to fit a zero is hygiene. What the
combination buys is that a practitioner can act on the finding without first
believing it. The checker tells an analyst that a forecast is internally
inconsistent using only what the analyst already has, and the closed model gives somewhere to go that
costs nothing in accuracy. A criticism that arrives with a remedy priced at zero is
adopted; one that arrives as a warning is filed.

\paragraph{A design-choice defence.}
We are not arguing for a link function. The repair that matters is dropping the
fitted aggregate, which requires no link at all, and the compositional form is what
that repair turns out to be equivalent to rather than an additional commitment. This
matters for adoption: an analyst who distrusts the softmax can decline it and still
take the entire correction, because Equation~\ref{eq:derived} is arithmetic. The
transferable knowledge is the consistency condition and the identification of the
aggregate as the term that must give way, not any particular parameterisation.

\paragraph{Why the defect survives ordinary quality control.}
Our uncoupled fits passed every gate we had. The residuals were small, the exponents
were well separated, the intervals were narrow, and the category rates declined
monotonically. None of those checks is sensitive to the violated identity, because
each examines one category at a time and the identity is a statement about the
categories jointly. This is the practical lesson for anyone maintaining an
evaluation pipeline: goodness-of-fit is computed within a series, and the constraint
lives across them.

\paragraph{Design-knowledge contribution.}
Stated at the level at which it transfers: when a quantity is decomposed and the
components are modelled separately, the identity that made the decomposition
meaningful becomes an unenforced constraint, and extrapolation is where the absence
of enforcement becomes visible. The remedy is to move the constraint into the
parameterisation rather than to check it after the fact, and the diagnostic value of
checking it after the fact is that it tells you how far the unconstrained model
could be trusted. This applies wherever errors are decomposed by capability,
benchmark subset, severity, customer segment, or failure type, and the test costs
one function call.

% ====================================================================
\section{Implications for Individuals, Organisations, and Society}
\label{sec:implications}

\paragraph{For the people affected by the systems these forecasts govern.}
The forecasts in question decide which defects get fixed and which are left to be
cured by the next model. That decision is felt by whoever is on the receiving end of
the defect that was left: the applicant whose document is misread, the patient whose
triage note is misfiled, the small trader whose claim is misclassified. When a
category of failure is forecast to disappear with scale, mitigation for it is
deferred, and the deferral lasts as long as the forecast is believed. An
inconsistent forecast systematically favours deferral, because the categories it
inflates are the shallow ones, which are exactly the ones described as persistent.
The people who carry the cost of an over-optimistic roadmap are rarely the people
who wrote it, and they have no way to contest an analysis they cannot see. A test
that runs on published coefficients gives an advocate, a regulator, or an
ombudsman something concrete to ask for.

\paragraph{For managers and organisations.}
The practical reframing is from a number to be received into a claim to be governed.
A supplier asserting that a failure class is on a steep curve is making a forecast
with an admissibility property, and that property can be checked before the contract
is signed rather than after the upgrade disappoints. We would put the question
directly into evaluation practice: ask for the fitted coefficients, run the check,
and require that the forecast scale lie below the breakdown scale. This is cheap
enough to be routine, it needs no access to the supplier's data or weights, and it
is auditable, so the reasoning survives the departure of whoever did it. It also
disciplines internal planning, where the same pattern is used to decide whether to
staff a mitigation team or wait for the next generation. Waiting is a decision with
a cost, and it should not rest on arithmetic nobody checked.

\paragraph{Strategic implications.}
Treat forecast admissibility as an ownership question. An organisation that can
audit capability claims in-house is not dependent on its suppliers' analysts for the
most consequential input into its own roadmap, and it is not exposed when a vendor
relationship ends. The artifact is small, has one dependency, and requires no
proprietary data, so this is genuinely available to organisations that could not
build an evaluation science function. We also read it as an accountability asset
rather than an overhead. A firm that can show it tested the admissibility of the
forecasts underlying a deployment decision is in a materially stronger position with
auditors and regulators than one that can show only that it believed a vendor, and
the record is created as a by-product of doing the analysis properly.

\paragraph{For society and digital equity.}
The ability to check a scaling claim should not depend on having a research team. A
diagnostic that consumes only published coefficients is runnable by a civil-society
organisation, a journalist, or a procurement officer in a small public body, on
ordinary hardware. We also record what the artifact deliberately does not do: it
does not rank models, score vendors, or certify anything, but answers the single
question of whether a forecast is entitled to the numbers it reports.

% ====================================================================
\section{Limitations}
\label{sec:limits}

\begin{description}[nosep,leftmargin=1.5em]
\item[L1. One instrument domain.] Both families are chess transformers, chosen
because a rules engine gives an exact and unambiguous error partition. The
admissibility argument is arithmetic and does not depend on the domain, but the
empirical demonstration does.
\item[L2. Consistency is not validation.] The derived total is arithmetically
possible, which the fitted aggregate is not. That does not make it observed. Every
scale beyond 39.9M is extrapolation, and consistency constrains the
form of a forecast without supplying evidence for it.
\item[L3. Three checkpoints in the independent family.] Its coefficients are
two-parameter fits to three points with one category dropped, so the replication is
reported as directional and no figure from it should be read as an estimate.
\item[L4. The invalidation claim is narrow.] What an inconsistent pair voids are
conclusions depending on its out-of-range or share-space behaviour, since that
behaviour is fixed by the misspecification rather than by data. Conclusions resting
on in-range residuals are unaffected, and we make no broader claim.
\item[L5. Depth and width are confounded.] Both ladders grow depth and width
together, so an exponent in parameters isolates neither.
\item[L6. Partitions only.] The condition is derived for categories that partition a
total. Overlapping categories are common in practice and need a different bound,
which we have not derived.
\item[L7. The audit is an existence check, not a survey.] We establish that speech
recognition already derives its total and that our own companion study did not. We
have not estimated how often the failing configuration occurs in the scaling
literature, which is the question a practitioner would most like answered and which
needs the survey of FR6.
\item[L8. The tolerance is a convention.] Ten percent is a defensible default and
not a derived quantity; a different tolerance moves the reported breakdown scale
though not the verdict.
\end{description}

% ====================================================================
\section{Avenues of Future Research}
\label{sec:future}

\begin{description}[nosep,leftmargin=1.5em]
\item[FR1] Apply the diagnostic in a domain whose error partition comes from a
non-game oracle, such as program execution traces or spreadsheet recalculation,
where a checker can say not only that a step is wrong but which constraint it
violated, to establish that the empirical demonstration is not an artefact of chess.
\item[FR2] Compare admissible links, including the additive-logratio and
multinomial-logit families, on their out-of-range behaviour, so that an analyst
choosing among them has evidence rather than only the guarantee they share.
\item[FR3] Extend the independent replication to a family with more checkpoints, so
that admissibility can be tested against exponents estimated with usable precision
rather than only in direction.
\item[FR4] Characterise precisely which inferences an inconsistent pair voids, by
constructing cases where a conclusion depends on in-range residuals alone and
confirming it survives, so that the narrow claim of L4 can be stated as a criterion
rather than as a caution.
\item[FR5] Separate depth from width along a ladder that varies them independently,
and determine whether the divergence rate is a property of parameter count or of
one of its components.
\item[FR6] Derive the corresponding bound for overlapping categories, which would
extend the check to severity scales and multi-label failure taxonomies as used in
industrial evaluation.
\item[FR7] Conduct a survey of published per-category scaling forecasts, applying
the diagnostic to their reported coefficients, to estimate how frequently forecasts
are read beyond their own breakdown scale.
\item[FR8] Study the tolerance itself, relating it to the sampling error of the
underlying rate estimates so that the breakdown scale can be reported with a
principled rather than a conventional threshold.
\end{description}

% ====================================================================
\section{Conclusion}
\label{sec:conclusion}

Decomposing a system's errors and fitting a scaling curve to each category is a
reasonable thing to want to do, and the resulting forecast decides where money and
attention go. What cannot also be done is to fit the aggregate error rate as a curve
of its own, because a sum of power laws with differing exponents is not a power law,
and the two reports then contradict each other by a margin that grows without bound.
On our measurements the contradiction passed ten percent beyond
174M parameters [57M, 635M],
and we had read a forecast at one billion. The repair is to derive the total from
the parts, which is arithmetic rather than methodology and which speech recognition
adopted long ago. Its price is the aggregate exponent, a number that does not
survive contact with the constraint: on our data its effective slope moves
0.125 between the bottom of our range and one billion parameters.

The finding we would most want carried forward is the one we obtained by being
wrong. An inconsistent pair does not merely produce forecasts that are too
confident. Where its out-of-range behaviour is fixed by the misspecification rather
than by data, that behaviour cannot be evidence for or against anything, so
conclusions drawn from it in either direction are void. We found this because a
hypothesis we had rejected turned out to have been rejected by nothing. The people this ultimately matters for are not the analysts.
They are the users of systems whose known defects were left unrepaired because a
forecast said scale would take care of them, and who had no way to ask whether the
arithmetic held.

% ====================================================================
\section*{Declarations}

\paragraph{Disclosure statement.}

Removed for double-anonymous peer review.


\paragraph{Funding.}

Removed for double-anonymous peer review.


\paragraph{Data availability statement.}

The code, the measurement artefacts, and the generated manuscript sources will be
made available upon acceptance.


\paragraph{AI usage declaration.}
During the preparation of this manuscript the authors used a generative AI tool to
assist with language editing and refinement of the text. The authors reviewed and
edited all output and take full responsibility for the content of the paper.
Generative AI was not used to generate or analyse the study's data, which were
produced entirely by the evaluation pipeline described in the paper.

\bibliographystyle{apalike}
\bibliography{refs}

\end{document}
